\newpage
\chapter{Instrumentación y adquisición de señales}
\vspace{3\baselineskip}

Este capítulo detalla los fundamentos teóricos de los sistemas de adquisición de datos en el contexto de ensayos de alta tensión. Se aborda la transición del dominio analógico al digital para la captura de ondas de impulso tipo rayo (\SI[parse-numbers=false]{1,2/50}{\micro\second}) empleando un osciloscopio de almacenamiento digital.

A lo largo de las siguientes secciones, se analizan los conceptos básicos necesarios para entender el funcionamiento de un osciloscopio.

\section{Osciloscopio de almacenamiento digital}

El osciloscopio de almacenamiento digital (DSO, \emph{Digital Storage Oscilloscope}), es un instrumento de medición electrónico, el cual se encarga de convertir la señal analógica a un formato digital, permitiendo su almacenamiento, visualización y posterior análisis.

Su funcionamiento consiste en que la señal ingresa por un circuito analógico de acondicionamiento de entrada, para dirigirse al bloque del convertidor analógico-digital (ADC, \emph{Analog-Digital Converter}) que muestrea y cuantiza la señal a una tasa definida por la base de tiempo. Luego, se almacena la verisón digitalizada en la memoria del osciloscopio, para finalmente presentarla en pantalla.

A continuación, se describen sus parámetros más relevantes, que deben ser considerados al momento de seleccionar el osciloscopio para capturar impulsos de alta tensión.

\textbf{Muestreo de señales:} Consiste en tomar valores de una señal analógica continua a intervalos regulares de tiempo para transformarla en una señal discreta de valores continuos (infinitos valores posibles). La conversión analógica-digital se basa en el teorema de muestreo, el cual establece que la frecuencia de muestreo $(f_s)$ debe ser mayor al doble de la componente de máxima frecuencia $(f_{max})$ presente en la señal, ya que si la velocidad de muestreo es insuficiente, pueden aparecer el fenómeno destructivo del \enquote{aliasing} que muestra erróneamente componentes de frecuencias menores que la señal real, denominadas “Alias”. Matemáticamente, se define:
\begin{equation}
    f_s > 2 f_{max}
\end{equation}

\textbf{Resolución vertical:} La resolución del eje vertical de un osciloscopio está vinculada con la resolución del ADC que tiene internamente, la cual está dada por la cantidad de bits ($N$) del mismo. Consierando que la tensión de entrada varia entre cero y un valor máximo $(V_{max})$, matemáticamente, puede expresarse:
\begin{equation}
    V_{res} = \dfrac{V_{max}}{2^N}
\end{equation}

\textbf{Cuantización de señales:} Consiste de convertir el valor de amplitud continuo en un valor discreto,  dentro de una escala predeterminada. Este proceso introduce un error numérico inherente denominado \textbf{error de cuantización}, ocasionado por la aproximación del valor real al nivel discreto más cercano. El tamaño del paso de cuantización o resolución del LSB (Bit Menos Significativo, por sus siglas en inglés) está dado por:
\begin{equation}
    q = \dfrac{V_{FSR}}{2^N}
\end{equation}
Donde $V_{FSR}$ representa el rango dinámico de escala completa (FSR) (relación entre el valor máximo y el valor mínimo de la señal que se puede capturar o representar), y $N$ es la resolución en bits del convertidor.

\textbf{Longitud de memoria}: El osciloscopio posee una memoria finita ($N_{record}$) para almacenar datos, que junto con la frecuencia de muestreo, determina la ventana de tiempo ($T_{acq}$) que puede registrar datos, a una frecuencia de muestreo ($f_s$) dada.
\begin{equation}
    T_{acq} = \dfrac{N_{record}}{f_s}
\end{equation}

\textbf{Ancho de banda:} Es un parámetro que se mide en \unit{MHz} o \unit{GHz}, e indica la frecuencia máxima a la que puede responder con precisión. Para garantizar una representación fiel de la señal, el ancho de banda debe ser superior a la frecuencia máxima de la señal que se desea medir.

\textbf{Velocidad de muestreo:} Es una medida de la resolución en el tiempo, que se mide en \unit{MSa/s} o \unit{GSa/s}, e indica la cantidad de muestras por unidad de tiempo (normalmente en segundos) que captura el instrumento, por lo tanto, define implícitamente la separación horizontal entre dos muestras consecutivas.

\section{Estándares de comunicación de instrumentos}

En el diseño e implementación de sistemas de instrumentación virtual, la transferencia de los datos desde la memoria del osciloscopio hacia la computadora requiere de interfaces físicas estandarizadas, y protocolos lógicos de control.

Actualmente, la comunicación entre instrumentos de medida y sistemas de adquisición se rige por los siguientes estándares físicos:

\begin{itemize}
    \item \textbf{GPIB (General Purpose Interface Bus / IEEE-488):} Es un estándar de bus de datos digital de corto rango diseñado específicamente para interconectar y controlar dispositivos de prueba y medida, tales como osciloscopios, multímetros y generadores de señal, mediante un controlador central.
    \item \textbf{Comunicaciones Serie (RS-232 y USB):} Los datos resultantes del muestreo y la digitalización suelen transmitirse a la computadora personal empleando el puerto USB o el puerto Serie (utilizando el protocolo RS-232). Estas interfaces son de adopción masiva debido a la arquitectura abierta de las computadoras modernas.
    \item \textbf{LXI (LAN eXtensions for Instrumentation) / Ethernet:} Permite la integración del instrumento en redes de área local, facilitando tasas de transferencia de datos significativamente mayores y posibilitando el control y monitoreo remoto a largas distancias.
\end{itemize}

\subsection{Protocolo VISA}

La diversidad de buses físicos (GPIB, RS-232, USB, Ethernet) presenta un desafío a nivel de desarrollo de software, ya que tradicionalmente cada puerto requería llamadas al sistema y controladores (\emph{drivers}) de bajo nivel completamente distintos. Para mitigar este problema, la industria adoptó la Arquitectura de Software para Instrumentos Virtuales (VISA, \emph{Virtual Instrument Software Architecture}), que funciona como una capa de abstracción de Entrada/Salida, cuyo objetivo es proporcionar una interfaz de software, que encapsula y oculta los detalles de implementación del hardware subyacente. Lo que brinda las siguientes ventajas en el desarrollo de este trabajo:

\begin{itemize}
    \item \textbf{Independencia del bus de hardware:} VISA permite que el software de control envíe comandos al osciloscopio sin necesidad de especificar cómo se gestionan eléctricamente los bits. Si el diseño del sistema migra de una conexión USB a una LAN, el código fuente de control de nivel superior no requiere modificaciones sustanciales.
    \item \textbf{Integración con comandos SCPI:} VISA actúa como el canal de transporte para el protocolo de Comandos Estándar para Instrumentos Programables (SCPI, \emph{Standard Commands for Programmable Instruments}), que es un lenguaje basado en cadenas de texto ASCII estandarizadas (por ejemplo, \texttt{TIMebase:SCALe?}) y los instrumentos interpretan para ajustar sus parámetros como escala de tiempo, resolución vertical, o para entregar el contenido de su memoria hacia la computadora.
    \item \textbf{Gestión de recursos:} La arquitectura asigna descriptores unificados a cada instrumento detectado, centralizando la inicialización, la lectura de las muestras adquiridas, la escritura de comandos de configuración y el cierre de las sesiones de comunicación.
\end{itemize}

Los descriptores de recursos VISA por lo general, tienen la siguiente gramática estructural:
\begin{align*}
    \texttt{Interfaz[Índice]::Dirección\_del\_Dispositivo::[Recurso]}
\end{align*}

\begin{itemize}
    \item \textbf{Interfaz}: Indica el tipo de controlador del puerto.
    \item \textbf{Índice}: Permite diferenciar numéricamente (por ejemplo, el 0 en USB0) cuando existen múltiples tarjetas controladoras del mismo tipo instaladas en la misma computadora (por ejemplo, dos placas PCI para GPIB distintas). Si se omite, VISA asume el valor cero por defecto.
    \item \textbf{Dirección del Dispositivo}: Identificador o ruta simbólica asignada para ubicar de forma unívoca el dispositivo.
    \item \textbf{Recurso}:
    \begin{itemize}
        \item \textbf{\texttt{INSTR}}: La sesión se gestiona utilizando los mecanismos de VISA completos, permitiendo por ejemplo, el uso de comandos de aserción de eventos y de lectura de estado.
        \item \textbf{\texttt{SOCKET}}: VISA actuará simplemente como un puente datos de bajo nivel sin enviar comandos de control propios del protocolo de instrumentación.
    \end{itemize}
\end{itemize}

A continuación, se presentan ejemplos de descriptores de recursos VISA para cada tipo de interfaz física comúnmente utilizada en la adquisición de señales:

\begin{table}[H]
    \centering
    \caption{Ejemplo de descriptor para ASRL (Asynchronous Serial Line)}
    \begin{tabular}{@{}ll@{}}
        \toprule
        \multicolumn{2}{c}{\textbf{\texttt{ASRL1::INSTR}}} \\ \midrule
        \texttt{ASRL1} & Controlador lógico RS-232 o UART.                  \\
                       & Mapeado por el sistema al puerto físico (ej. COM1) \\ \bottomrule
    \end{tabular}
\end{table}

\begin{table}[H]
    \centering
    \caption{Ejemplo de descriptor para USB (Universal Serial Bus)}
    \begin{tabular}{@{}ll@{}}
        \toprule
        \multicolumn{2}{c}{\textbf{\texttt{USB0::0x0957::0x17A6::MY503::INSTR}}} \\ \midrule
        \texttt{USB0}   & Controlador lógico USB. \\
        \texttt{0x0957} & ID del fabricante.      \\
        \texttt{0x17A6} & ID del producto.        \\
        \texttt{MY503}  & Número de serie único.  \\ \bottomrule
    \end{tabular}
\end{table}